\documentclass[UTF8,zihao=-4,a4paper,fontset=fandol]{ctexart}
\usepackage{cumcm-paper}
\usepackage{makecell}
\usepackage{multirow}

\title{电商物流网络货量预测、关停分流与网络优化模型}
\author{}
\date{}

\begin{document}

\maketitle

\begin{cumcmabstract}
\hspace{2em}本文围绕电商物流网络的货量预测、场地关停分流和网络优化展开研究，建立异常修正预测、分层分流优化、重要性评价和新增场地设计模型，在统一能力口径下给出各问结果。

对于问题一，本文补全线路日度面板，修正短期异常和长期停运式异常，再用 LightGBM 递推预测并进行总量校准。回测得到 MAE 为 671.524、RMSE 为 4535.005、sMAPE 为 0.213583；问题一的答案为：2023 年 1 月全网预测总货量为 53,200,743 件，DC14->DC10、DC20->DC35、DC25->DC62 的月预测货量分别为 784,382、255、255,633 件。

对于问题二，本文在 DC5 关停且不新开线路条件下，建立按日求解的混合整数线性规划模型，依次优化未正常流转量、变化线路数和负荷压力，并比较三种变化线路放宽方案。问题二的答案为：采用每日最多多使用 2 条变化线路的折中方案，DC5 关停影响货量 3,214,028 件，正常分流 3,190,249 件，未正常流转 23,779 件，平均每日变化线路数为 44.741935 条。

对于问题三，本文在 DC9 关停且允许动态新开线路条件下，构造原有线路和新开线路候选集，比较最少变化 MILP、新开优先 MILP 和流启发式均衡三类方案。三类方案正常流转率相同，方案 A 的平均每日总变化线路数最低；问题三的答案为：采用方案 A，DC9 受影响货量 13,700,966 件，正常流转 13,688,493 件，未正常流转 12,473 件，正常流转率为 0.999090，平均每日总变化线路数为 37.161290 条。

对于问题四，本文将重要性评价与负荷风险分开处理，并结合核心场地、关键上下游和超载路径设计新增场地。问题四的答案为：场地重要性前五名为 DC14、DC9、DC10、DC3、DC4，线路重要性前五名为 DC14->DC9、DC25->DC9、DC14->DC3、DC9->DC3、DC9->DC14；新增场地 DC\_NEW 连接 14 个既有场地，新开 16 条有向线路，处理能力为 800,000 件/日。
\end{cumcmabstract}

\keywords{物流网络；LightGBM；混合整数线性规划；场地关停；货量分流；网络鲁棒性}

\CumcmCompactToc

\CumcmBodyStart

\section{问题重述}



电商物流网络由多个物流场地及其之间的有向运输线路构成。历史数据给出了 2021-01-01 至 2022-12-31 期间各线路的每日货量。要求首先预测 2023-01-01 至 2023-01-31 期间每条线路的每日货量，并重点给出 DC14->DC10、DC20->DC35、DC25->DC62 三条线路的预测结果。



在预测结果基础上，需要进一步分析场地关停后的网络调运能力。问题二假设 DC5 自 2023-01-01 起关停，且不允许新开线路，要求将所有与 DC5 相关的货量分配至其他已有线路，使包裹尽可能正常流转，同时使关停前后货量发生变化的线路数尽可能少，并保持线路负荷尽可能均衡。问题三进一步假设 DC9 关停，但允许动态新开线路，需要给出每日可执行的线路调整方案，并按正常流转、变化线路数和负荷均衡三个标准评价方案。问题四要求评价物流场地和线路的重要性，并在网络中新增一个物流场地及若干线路，以提升网络稳定性和抗风险能力。



上述问题之间具有明显递进关系：问题一提供全网正常运行状态下的预测基准；问题二和问题三在该基准上研究不同关停约束下的分流策略；问题四则根据预测货量、网络结构和关停影响识别关键节点与线路，进一步提出结构性优化方案。



\section{问题分析}



问题一的核心是得到可信的“正常运行基准”。历史数据中存在大量空白线路日期组合，也存在部分线路在某一阶段异常停运或货量异常下降的现象。若直接将这些异常作为正常历史规律学习，会导致后续关停问题中的基准需求被低估。例如 DC9 在 2022 年出现长期停运式异常，但问题三又假设 DC9 从 2023-01-01 起关停，这意味着若没有该关停，DC9 应存在正常货量。因此问题一必须先修正异常，再进行预测。



问题二和问题三都属于受限网络流重分配问题。二者差别在于，问题二不允许新开线路，分流只能在已有线路中完成；问题三允许动态新开线路，因此候选分配空间更大，但也必须控制网络变化幅度。两问的目标并非单一最小成本，而是具有优先级的多目标优化：正常流转优先于线路变化，线路变化优先于负荷改善。



问题四前半部分要求识别重要场地和重要线路，评价对象是“关停或异常后会对网络造成多大影响”，因此负荷率不应进入重要性得分。负荷率反映的是运行压力或超载风险，可用于后半部分新增场地与线路设计，但不能等同于重要性。后半部分新增场地设计则需要同时考虑关键节点保护、关键上下游覆盖、现有超载路径缓解和新增线路成本。



\section{基本假设与符号说明}



\subsection{基本假设}



\begin{enumerate}

\item 原始数据中某线路在某日无记录，表示该线路当日无货量流转，补为 0。

\item 线路运输能力取该线路历史日货量最大值，场地处理能力取该场地历史每日入货量与出货量之和的最大值。

\item 对短期异常和长期停运式异常，若其明显不代表正常需求下降，则先进行历史基准修正，再用于预测和后续关停优化。

\item 若预测基准本身已经超过历史能力，不直接修改能力上限，也不舍弃预测货量，而是将其记录为基准超载状态；后续分流时不再向无剩余能力的线路或场地追加货量。

\item 问题二中不允许新开线路；问题三中允许每日动态新开线路，且新线路能力上限不超过历史已有线路最大能力 505,833。

\item 问题四新增场地的线路方向根据相连场地的流向角色确定，不默认全部双向。

\end{enumerate}



\subsection{符号说明}




\begin{table}[H]
  \centering
  \caption{主要符号说明}
  \label{tab:auto-1}
  \begin{tabular*}{\textwidth}{@{}L{0.50\textwidth}@{\extracolsep{\fill}}L{0.42\textwidth}@{}}
    \toprule
符号 & 含义 \\
\midrule
$e=(i,j)$ & 从场地 $i$ 到场地 $j$ 的有向线路 \\
$y_{e,t}$ & 线路 $e$ 在日期 $t$ 的实际货量 \\
$\hat y_{e,t}$ & 线路 $e$ 在日期 $t$ 的预测货量 \\
$C_e$ & 线路 $e$ 的运输能力 \\
$S_k$ & 场地 $k$ 的处理能力 \\
$B_{e,t}$ & 关停或调整前线路 $e$ 在日期 $t$ 的基准预测货量 \\
$P_{k,t}$ & 场地 $k$ 在日期 $t$ 的处理量 \\
$x_{r,e,t}$ & 将受影响需求 $r$ 分配至候选线路 $e$ 的货量 \\
$u_{r,t}$ & 受影响需求 $r$ 未能正常流转的货量 \\
$z_{e,t}$ & 线路 $e$ 在日期 $t$ 是否发生承接变化 \\
$L_t$ & 日期 $t$ 的线路负荷均衡辅助变量 \\
    \bottomrule
  \end{tabular*}
\end{table}




\section{问题一：异常修正的全网货量预测模型}



\subsection{数据预处理}



首先将历史记录扩展为完整的“线路-日期”面板。设全网共有 1049 条有向线路，历史日期为 730 天，则完整面板包含



\begin{equation}
1049 \times 730 = 765770
\end{equation}



条样本。若原始附件中某条线路在某日无记录，则补为 0。该处理可以使模型同时学习活跃日期和非活跃日期，避免只在有货量样本上训练而系统性高估稀疏线路。



其次识别并修正异常数据。短期异常按月度局部中位数识别：若某线路某月货量显著低于前后邻近月份的正常水平，且前后月份未同步下降，则认为该月存在短期异常，用邻近月份日度形态进行修正。长期停运式异常按年度总量识别：若某线路 2021 年货量较高，但 2022 年货量不足 2021 年的 5\%，则认为其 2022 年可能处于异常停运状态，不能直接作为正常运营样本。该类线路使用 2021 年正常数据并结合全网稳定线路增长系数恢复 2022 年基准。



\subsection{特征构造}



对线路 $e$ 和日期 $t$，构造特征向量



\begin{equation}
\mathbf{x}_{e,t}=(\mathbf{x}^{line}_e,\mathbf{x}^{date}_t,\mathbf{x}^{festival}_t,\mathbf{x}^{lag}_{e,t},\mathbf{x}^{roll}_{e,t},\mathbf{x}^{stat}_e).
\end{equation}



其中，线路特征包括线路编号、起点编号和终点编号；日期特征包括月份、日、星期、是否周末、是否月初和是否月末；春节特征使用春节相对日期刻画。设 $s(t)$ 为日期 $t$ 所在年份的春节日期，则春节相对日期为



\begin{equation}
d_t=t-s(t).
\end{equation}



滞后特征选取 $1,2,3,7,14,28,365$ 日滞后货量，以刻画短期惯性、周周期和去年同期水平。滚动特征包括 7、14、28、56 日滚动均值、滚动最大值以及近 28 日非零比例。静态统计特征包括线路历史均值、中位数、最大值、非零比例，以及起点和终点场地的历史统计量。



\subsection{LightGBM 预测模型}



由于线路货量非负且不同线路量级差异明显，先对目标变量作对数变换：



\begin{equation}
q_{e,t}=\log(1+y_{e,t}).
\end{equation}



建立 LightGBM 加性回归树模型：



\begin{equation}
\hat q_{e,t}=F(\mathbf{x}_{e,t})=\sum_{m=1}^{M} f_m(\mathbf{x}_{e,t}),
\end{equation}



再反变换得到原货量空间预测值：



\begin{equation}
\hat y^{raw}_{e,t}=\max\{0,\exp(\hat q_{e,t})-1\}.
\end{equation}



预测 2023 年 1 月时，滞后特征依赖前序预测结果，因此采用逐日递推预测：先预测 2023-01-01，将其写入历史序列，再生成 2023-01-02 的滞后和滚动特征，依次递推至 2023-01-31。



\subsection{总量校准}



异常修正后，根据 2021 年和 2022 年春节相对窗口 $[-21,9]$ 的全网货量构造 2023 年 1 月目标总量：



\begin{equation}
Q^*=0.35Q_{2021}+0.65Q_{2022}.
\end{equation}



修正后有




\begin{table}[H]
  \centering
  \caption{问题一总量校准指标}
  \label{tab:auto-2}
  \begin{tabular*}{\textwidth}{@{}L{0.50\textwidth}@{\extracolsep{\fill}}L{0.42\textwidth}@{}}
    \toprule
指标 & 数值 \\
\midrule
2021 年春节窗口总量 $Q_{2021}$ & 37,284,654 \\
2022 年春节窗口总量 $Q_{2022}$ & 61,770,945 \\
2023 年 1 月目标总量 $Q^*$ & 53,200,743 \\
模型原始预测总量 & 57,425,198 \\
校准系数 $\alpha$ & 0.926436 \\
    \bottomrule
  \end{tabular*}
\end{table}




最终预测值为



\begin{equation}
\hat y_{e,t}=round(\alpha \hat y^{raw}_{e,t}).
\end{equation}



该校准保留了 LightGBM 给出的线路间相对结构，同时使全网总量与春节相对窗口规律一致。




\begin{figure}[H]
  \centering
  \includegraphics[width=0.95\textwidth]{\detokenize{../../q1_revised/figures/daily_total_before_after.png}}
  \caption{问题一全网日预测总量修正前后对比}
\end{figure}




\subsection{预测结果}



修正后模型在 2022 年 1 月正常基准上的回测指标为：




\begin{table}[H]
  \centering
  \caption{问题一预测模型回测误差}
  \label{tab:auto-3}
  \begin{tabular}{L{0.28\textwidth}L{0.28\textwidth}L{0.34\textwidth}}
    \toprule
MAE & RMSE & sMAPE \\
\midrule
671.524 & 4535.005 & 0.213583 \\
    \bottomrule
  \end{tabular}
\end{table}




2023 年 1 月全网预测总量为 53,200,743 件。三条指定线路的月预测总量如下：




\begin{table}[H]
  \centering
  \caption{三条指定线路月预测货量}
  \label{tab:auto-4}
  \begin{tabular*}{\textwidth}{@{}L{0.50\textwidth}@{\extracolsep{\fill}}L{0.42\textwidth}@{}}
    \toprule
指定线路 & 2023 年 1 月预测总量 \\
\midrule
DC14->DC10 & 784,382 \\
DC20->DC35 & 255 \\
DC25->DC62 & 255,633 \\
    \bottomrule
  \end{tabular*}
\end{table}




修正前后关键变化包括：DC9 相关货量由 0 恢复至 13,700,966 件，DC3 相关货量由 1,664 件恢复至 5,108,060 件。该处理使后续 DC9 关停问题具有合理的正常运行基准。




\begin{figure}[H]
  \centering
  \includegraphics[width=0.95\textwidth]{\detokenize{../../q1_revised/figures/key_sites_before_after.png}}
  \caption{问题一关键场地预测货量修正前后对比}
\end{figure}




\section{问题二：DC5 关停下的已有线路分流模型}



\subsection{受影响需求定义}



DC5 关停后，所有与 DC5 相连且预测货量为正的线路均需重新分配。若原需求为 $i\to DC5$，则候选替代线路为 $i\to k$，即保持原起点不变，替换终点；若原需求为 $DC5\to j$，则候选替代线路为 $k\to j$，即保持原终点不变，替换起点。问题二不允许新开线路，因此候选线路必须为历史网络中已存在且非 DC5 相关的线路。



基于问题一修正预测结果，DC5 相关受影响总货量为 3,214,028 件，其中入 DC5 货量为 3,161,350 件，出 DC5 货量为 52,678 件。



\subsection{MILP 决策变量}



对每日 $t$、每条受影响需求 $r$ 和候选替代线路 $e$，定义：




\begin{table}[H]
  \centering
  \caption{问题二 MILP 决策变量}
  \label{tab:auto-5}
  \begin{tabular*}{\textwidth}{@{}L{0.50\textwidth}@{\extracolsep{\fill}}L{0.42\textwidth}@{}}
    \toprule
变量 & 含义 \\
\midrule
$x_{r,e,t}\ge 0$ & 分配至候选线路 $e$ 的货量，取整数 \\
$u_{r,t}\ge 0$ & 需求 $r$ 未正常流转的货量，取整数 \\
$z_{e,t}\in\{0,1\}$ & 线路 $e$ 是否因承接分流而发生货量变化 \\
$L_t\ge 0$ & 当日线路负荷辅助变量 \\
    \bottomrule
  \end{tabular*}
\end{table}




\subsection{约束条件}



需求守恒约束为



\begin{equation}
\sum_{e\in A(r)}x_{r,e,t}+u_{r,t}=d_{r,t},
\end{equation}



其中 $A(r)$ 为需求 $r$ 的候选线路集合，$d_{r,t}$ 为原 DC5 相关需求量。



线路剩余能力约束为



\begin{equation}
\sum_r x_{r,e,t}\le \max(C_e-B_{e,t},0).
\end{equation}



若某线路在基准预测中已经超出历史能力，则其剩余能力记为 0，不再承接新增分流。



场地处理能力约束为



\begin{equation}
\sum_{e\in \delta(k)} a^{new}_{e,k,t} \le \max(S_k-P_{k,t},0),
\end{equation}



其中 $P_{k,t}$ 为场地 $k$ 在基准预测下的当日处理量。基准已超处理能力的场地也不再承接新增分流。



线路变化指示约束为



\begin{equation}
\sum_r x_{r,e,t}\le M_e z_{e,t}.
\end{equation}



负荷约束为



\begin{equation}
\frac{B_{e,t}+\sum_r x_{r,e,t}}{C_e}\le L_t.
\end{equation}



\subsection{分层优化目标}



题目目标存在明确优先级，因此采用三阶段逐日 MILP。每天先根据预测基准提取 DC5 相关受影响需求和各候选已有线路的剩余能力，再依次执行以下三个优化阶段：



第一阶段最小化未正常流转货量：



\begin{equation}
\min \sum_r u_{r,t}.
\end{equation}



第二阶段在第一阶段最优值固定的前提下，最小化变化线路数：



\begin{equation}
\min \sum_e z_{e,t}.
\end{equation}



第三阶段在未流转货量最小且变化线路数基本不变的前提下，降低线路负荷压力。设第二阶段得到的最少变化线路数为 $z_t^{min}$，放宽量为 $\delta$，则有



\begin{equation}
\min L_t,\quad \sum_e z_{e,t}\le z_t^{min}+\delta.
\end{equation}



本文取 $\delta=0,2,5$ 构造三种方案进行比较：$\delta=0$ 表示严格采用最少变化线路数；$\delta=2$ 表示在每日最多增加 2 条变化线路的范围内调整负荷；$\delta=5$ 表示给负荷优化更大的调整空间。这样可以检验在变化线路数基本不变的条件下，额外线路调整是否能换来更好的负荷表现，避免为了极小幅度改善负荷而打开大量线路。



\subsection{三种放宽方案对比}



三种方案均以第一阶段最小未正常流转量为硬约束，因此正常流转能力保持一致；差别在于第二阶段最少变化线路数附近保留多少负荷调整空间。结果如下：




\begin{table}[H]
  \centering
  \caption{问题二变化线路放宽方案对比}
  \label{tab:auto-6}
  \resizebox{\textwidth}{!}{%
  \begin{tabular}{llllllll}
    \toprule
方案 & 正常求解天数 & 受影响货量 & 未正常流转货量 & 平均每日变化线路数 & 最大每日变化线路数 & 最大线路负荷率 & 负荷率日均标准差 \\
\midrule
$+0$ 严格最少变化 & 31/31 & 3,214,028 & 23,779 & 42.741935 & 63 & 1.204408 & 0.102321 \\
$+2$ 小幅放宽 & 31/31 & 3,214,028 & 23,779 & 44.741935 & 65 & 1.204408 & 0.102743 \\
$+5$ 较大放宽 & 31/31 & 3,214,028 & 23,779 & 46.645161 & 67 & 1.204408 & 0.105081 \\
    \bottomrule
  \end{tabular}}
\end{table}




从正常流转看，三种方案 31 天均能求得最优解，未正常流转总量均为 23,779 件，说明第一优先级由网络剩余能力决定，变化线路放宽不会改变可流转上限。从变化线路数看，$+0$ 方案平均每日变化线路数最少，为 42.741935 条；$+2$ 方案平均每日增加约 2 条变化线路，仍保持在“变化基本不变”的范围内；$+5$ 方案变化线路数进一步增加，但未带来正常流转改善。



从负荷表现看，新版预测基准下最大线路负荷率均为 1.204408，说明最大负荷主要受基准中已有超载线路控制，而不是 DC5 分流本身造成。负荷率日均标准差在 $\delta=0,2,5$ 下分别为 0.102321、0.102743、0.105081，较大放宽并未形成稳定的均衡改善。因此本文不采用 $\delta=5$；$\delta=0$ 可作为严格少变化参考，$\delta=2$ 则作为推荐折中方案，用少量额外变化为负荷调整保留可执行余量，同时不显著扩大网络改动幅度。



\subsection{结果分析}



采用 $\delta=2$ 折中方案时，DC5 关停分流结果如下：




\begin{figure}[H]
  \centering
  \includegraphics[width=0.95\textwidth]{\detokenize{figures/table7_dc5_diversion_bar.png}}
  \caption{问题二 DC5 关停分流结果}
  \label{fig:dc5-diversion-bar}
\end{figure}




未正常流转货量只出现在 3 天：




\begin{table}[H]
  \centering
  \caption{问题二未正常流转日期与货量}
  \label{tab:auto-8}
  \begin{tabular*}{\textwidth}{@{}L{0.50\textwidth}@{\extracolsep{\fill}}L{0.42\textwidth}@{}}
    \toprule
日期 & 未正常流转货量 \\
\midrule
2023-01-01 & 4,718 \\
2023-01-13 & 16,196 \\
2023-01-16 & 2,865 \\
    \bottomrule
  \end{tabular*}
\end{table}




结果说明，在不新开线路的限制下，DC5 的绝大多数货量仍可通过已有网络完成调运。少量未正常流转主要来自预测基准恢复后网络负荷提高，以及部分场地或线路在基准状态下已经接近或超过历史能力。综合正常流转、变化线路数和负荷均衡表现，$+2$ 方案比 $\delta=0$ 多使用极少量变化线路以保留负荷调整余量，又比 $\delta=5$ 更能控制改动规模，因此作为问题二最终方案。



\section{问题三：DC9 关停下的动态新开线路分流模型}



\subsection{建模思路}



DC9 是问题三中被指定关停的场地。由于问题一修正后恢复了 DC9 的正常运行基准，DC9 相关受影响货量达到 13,700,966 件，其中入 DC9 货量为 12,413,163 件，出 DC9 货量为 1,287,803 件。若不允许新开线路，仅依靠原有线路难以获得足够灵活的承接空间，因此问题三允许动态新开线路。



本文仍按逐日优化处理。对原 $i\to DC9$ 需求，可分配至 $i\to k$；对原 $DC9\to j$ 需求，可分配至 $k\to j$。其中候选线路既包括已有线路，也包括当日允许新开的线路。新线路能力上限取历史已有线路最大能力 505,833。



\subsection{评价标准}



根据题意和前述讨论，问题三采用三个标准，且优先级如下：



\begin{enumerate}

\item 正常流转能力：看受影响货量中有多少能够被成功分流，核心指标为正常流转货量、未正常流转货量和正常流转率。

\item 线路变化幅度：看关停后网络调整是否尽量少，核心指标为关闭线路数、承接变化线路数、总变化线路数。

\item 承接线路负荷均衡：不再使用最大负荷率和平均负荷率作为主要比较口径，而使用承接线路负荷率的标准差和变异系数刻画均衡程度。

\end{enumerate}



\subsection{三类方案对比}



本文比较三类方案：




\begin{table}[H]
  \centering
  \caption{问题三三类方案建模思路对比}
  \label{tab:auto-9}
  \begin{tabular*}{\textwidth}{@{}L{0.50\textwidth}@{\extracolsep{\fill}}L{0.42\textwidth}@{}}
    \toprule
方案 & 核心思想 \\
\midrule
方案 A：最少变化 MILP & 优先保证正常流转，再尽量减少变化线路数，最后在少量变化余量内改善负荷 \\
方案 B：新开优先 MILP & 优先使用新开线路承接 DC9 货量，减少对原有线路的扰动 \\
方案 C：流启发式均衡 & 使用更快的启发式规则将货量分散至多条候选线路以改善负荷分布 \\
    \bottomrule
  \end{tabular*}
\end{table}




三类方案的标准口径对比如下：




\begin{table}[H]
  \centering
  \caption{问题三三类方案统一指标对比}
  \label{tab:auto-10}
  \resizebox{\textwidth}{!}{%
  \begin{tabular}{lrrr}
    \toprule
指标 & 方案 A 最少变化 MILP & 方案 B 新开优先 MILP & 方案 C 流启发式均衡 \\
\midrule
受影响总货量 & 13,700,966 & 13,700,966 & 13,700,966 \\
正常流转货量 & 13,688,493 & 13,688,493 & 13,688,493 \\
未正常流转货量 & 12,473 & 12,473 & 12,473 \\
正常流转率 & 0.999090 & 0.999090 & 0.999090 \\
平均每日关闭线路数 & 19.000000 & 19.000000 & 19.000000 \\
平均每日承接变化线路数 & 18.161290 & 20.709677 & 27.612903 \\
平均每日总变化线路数 & 37.161290 & 39.709677 & 46.612903 \\
最大每日总变化线路数 & 38 & 41 & 61 \\
承接线路负荷率日均标准差 & 0.252799 & 0.078135 & 0.149420 \\
承接线路负荷率日均变异系数 & 1.110176 & 1.641425 & 1.134901 \\
平均每日新开线路数 & 10.161290 & 20.451613 & 17.967742 \\
平均每日原有承接线路数 & 8.000000 & 0.258065 & 9.645161 \\
    \bottomrule
  \end{tabular}}
\end{table}




三个方案的正常流转能力完全相同，均仅在 2023-01-16 出现 12,473 件未正常流转，主要原因是当天 DC15 存在基准场地超载，模型不能再向无剩余能力的场地追加货量。在第一优先级相同的情况下，比较第二优先级，即变化线路数。方案 A 的平均每日总变化线路数为 37.161290，低于方案 B 的 39.709677 和方案 C 的 46.612903，因此更符合“变化尽可能少”的要求。



从负荷均衡看，方案 B 的负荷率标准差较低，但其几乎完全依赖新开线路，导致新开线路数量显著增加，且变异系数较高，说明承接线路之间的相对差异仍然明显。方案 C 的均衡性居中，但总变化线路数明显偏多。综合三项标准的优先级，本文选择方案 A 作为问题三最终方案。



\subsection{方案 A：受影响需求与候选线路}



方案 A 采用“最少变化 MILP”。其核心思路是：先把 DC9 关停导致的全部正货量线路转化为待分流需求，再在“起点或终点尽量保持不变”的原则下构造已有线路和新开线路候选集。这样既保留原物流流向的业务含义，又利用题目允许动态新开线路的条件扩大可行空间。



设日期 $t$ 下 DC9 相关受影响需求集合为 $R_t$。若原线路为 $i\to DC9$，则需求记为 $r=(i,DC9,t)$，候选承接线路为



\begin{equation}
A(r)=\{(i,k):k\ne i,\ k\ne DC9\}.
\end{equation}



若原线路为 $DC9\to j$，则需求记为 $r=(DC9,j,t)$，候选承接线路为



\begin{equation}
A(r)=\{(k,j):k\ne j,\ k\ne DC9\}.
\end{equation}



当候选线路 $e$ 已存在于历史网络中时，记为原有线路，能力取历史日货量最大值 $C_e$；当候选线路不存在时，允许作为当日新开线路，能力上限取全网已有线路能力最大值



\begin{equation}
C^{new}=505833.
\end{equation}



基于问题一修正预测结果，2023 年 1 月 DC9 相关受影响总货量为 13,700,966 件，其中入 DC9 货量为 12,413,163 件，出 DC9 货量为 1,287,803 件。



\subsection{方案 A：MILP 决策变量}



对每日 $t$、受影响需求 $r\in R_t$ 和候选承接线路 $e\in A(r)$，定义如下变量：




\begin{table}[H]
  \centering
  \caption{问题三方案 A MILP 决策变量}
  \label{tab:auto-11}
  \begin{tabular*}{\textwidth}{@{}L{0.50\textwidth}@{\extracolsep{\fill}}L{0.42\textwidth}@{}}
    \toprule
变量 & 含义 \\
\midrule
$x_{r,e,t}\ge 0$ & 将需求 $r$ 分配至候选线路 $e$ 的货量，取整数 \\
$u_{r,t}\ge 0$ & 需求 $r$ 未能正常流转的货量，取整数 \\
$z_{e,t}\in\{0,1\}$ & 候选线路 $e$ 是否因承接分流而发生货量变化 \\
$n_{e,t}\in\{0,1\}$ & 候选线路 $e$ 是否作为新开线路被启用 \\
$L_t\ge 0$ & 当日承接线路负荷率上界辅助变量 \\
    \bottomrule
  \end{tabular*}
\end{table}




其中 $z_{e,t}$ 用于计算承接变化线路数，$n_{e,t}$ 用于区分新开线路和原有承接线路。对原有线路有 $n_{e,t}=0$；对新开候选线路，若 $z_{e,t}=1$，则该线路计入当日新开线路。



\subsection{方案 A：约束条件}



需求守恒约束保证每条 DC9 相关需求要么被分流，要么计入未正常流转：



\begin{equation}
\sum_{e\in A(r)}x_{r,e,t}+u_{r,t}=d_{r,t},\quad r\in R_t.
\end{equation}



其中 $d_{r,t}$ 为原 DC9 相关线路在日期 $t$ 的预测货量。



对原有候选线路，新增分流不能超过其剩余运输能力。设 $B_{e,t}$ 为关停前线路 $e$ 的基准预测货量，则



\begin{equation}
\sum_{r:e\in A(r)}x_{r,e,t}\le \max(C_e-B_{e,t},0).
\end{equation}



若原有线路在基准预测中已超过历史能力，则剩余能力记为 0，不再承接新增分流。对新开候选线路，有



\begin{equation}
\sum_{r:e\in A(r)}x_{r,e,t}\le C^{new} n_{e,t}.
\end{equation}



场地处理能力约束用于防止分流后把压力转移到已经无剩余能力的场地。设 $P_{k,t}$ 为非 DC9 场地 $k$ 在基准预测下的当日处理量，$S_k$ 为其历史处理能力，则



\begin{equation}
\sum_{r,e:k\in e}x_{r,e,t}\le \max(S_k-P_{k,t},0),\quad k\ne DC9.
\end{equation}



线路变化指示约束为



\begin{equation}
\sum_{r:e\in A(r)}x_{r,e,t}\le M_e z_{e,t},
\end{equation}



其中 $M_e$ 可取该线路可用能力上界。若候选线路被分配正货量，则 $z_{e,t}=1$，该线路计入承接变化线路。



为在第三阶段改善负荷分布，定义候选承接线路的负荷率为



\begin{equation}
\rho_{e,t}=\frac{B_{e,t}+\sum_r x_{r,e,t}}{C_e},
\end{equation}



其中原有线路使用其历史能力 $C_e$，新开线路取 $C_e=C^{new}$ 且 $B_{e,t}=0$。



并用线性上界变量表示



\begin{equation}
\rho_{e,t}\le L_t.
\end{equation}



由于负荷率标准差含二次项，直接放入 MILP 会显著增加求解难度，因此模型求解阶段用 $L_t$ 作为线性化负荷均衡代理；最终方案评价时，再对实际承接线路负荷率计算标准差和变异系数。



\subsection{方案 A：分层优化目标}



题目要求首先保证正常流转，其次减少变化线路数，最后考虑负荷均衡。方案 A 因此采用三阶段逐日求解。



第一阶段最小化未正常流转货量：



\begin{equation}
U_t^*=\min \sum_{r\in R_t}u_{r,t}.
\end{equation}



该阶段决定当日最大可正常流转货量，优先级最高。



第二阶段在 $U_t^*$ 固定的前提下最小化承接变化线路数：



\begin{equation}
Z_t^*=\min \sum_e z_{e,t},\quad \sum_{r\in R_t}u_{r,t}=U_t^*.
\end{equation}



由于 DC9 关停导致的原 DC9 正货量线路必须关闭，关闭线路数在当天是固定值，所以最小化总变化线路数等价于最小化承接变化线路数。



第三阶段在未正常流转量最优且变化线路数基本不变的条件下改善负荷：



\begin{equation}
\min L_t,
\end{equation}



\begin{equation}
\sum_{r\in R_t}u_{r,t}=U_t^*,\quad
\sum_e z_{e,t}\le Z_t^*+2.
\end{equation}



其中 $+2$ 表示每日最多允许比最少变化方案多启用 2 条承接线路。该设置使模型不会为了追求局部负荷改善而大量新开线路，同时给负荷调整保留少量空间。



\subsection{方案 A：求解结果}



对 2023-01-01 至 2023-01-31 共 31 天分别建立并求解上述 MILP，31 个日模型均达到最优。方案 A 的结果如下：




\begin{table}[H]
  \centering
  \caption{问题三方案 A 分流结果}
  \label{tab:auto-12}
  \begin{tabular*}{\textwidth}{@{}L{0.50\textwidth}@{\extracolsep{\fill}}L{0.42\textwidth}@{}}
    \toprule
指标 & 数值 \\
\midrule
DC9 相关受影响总货量 & 13,700,966 \\
入 DC9 货量 & 12,413,163 \\
出 DC9 货量 & 1,287,803 \\
正常流转货量 & 13,688,493 \\
未正常流转货量 & 12,473 \\
正常流转率 & 0.999090 \\
平均每日关闭线路数 & 19.000000 \\
平均每日承接变化线路数 & 18.161290 \\
平均每日总变化线路数 & 37.161290 \\
最大每日总变化线路数 & 38 \\
平均每日新开线路数 & 10.161290 \\
平均每日原有承接线路数 & 8.000000 \\
承接线路负荷率日均标准差 & 0.252799 \\
承接线路负荷率日均变异系数 & 1.110176 \\
    \bottomrule
  \end{tabular*}
\end{table}




未正常流转货量只出现在 2023-01-16，当日未正常流转 12,473 件。核查出入方向可知，入 DC9 方向总货量 12,413,163 件，已分流 12,400,690 件，未正常流转 12,473 件；出 DC9 方向总货量 1,287,803 件，全部完成分流。该不可流转量主要由关停前基准预测中 DC15 场地已超过历史处理能力导致，属于基准超载背景下的剩余不可行量，不是 MILP 求解失败。



综合看，方案 A 在正常流转能力与其他方案相同的前提下，使平均每日总变化线路数保持最低。虽然方案 B 的负荷率标准差较小，但其平均每日新开线路数达到 20.451613 条，明显扩大网络调整规模；方案 C 的平均每日总变化线路数为 46.612903 条，也高于方案 A。因此，按照“正常流转优先、变化线路数其次、负荷均衡再次”的标准，方案 A 是问题三的最终推荐方案。



\section{问题四：场地和线路重要性评价及新增场地设计}



\subsection{重要性与负荷风险的区别}



场地或线路的重要性衡量的是其一旦失效会对网络货量和结构造成多大影响；负荷压力衡量的是其在当前预测状态下是否接近或超过能力上限。二者有关联，但不能混为一个指标。例如某节点负荷率很高，说明其短期运行风险大，但如果其货量规模小、网络连接少、可替代性强，则不一定是全网最重要节点。反之，某核心枢纽负荷率不高，但连接大量关键线路，关停后仍会造成巨大影响。



因此，本文在问题四前半部分的重要性评价中不引入负荷率；在后半部分新增场地和线路设计中，将负荷率、超载记录和超载货量作为风险改善指标。



\subsection{场地重要性评价模型}



记 $\hat y_{i,j,t}$ 为问题一得到的 2023 年 1 月线路 $i\to j$ 在日期 $t$ 的预测货量，$y_{i,j,t}^{his}$ 为历史货量。对场地 $k$，先计算其每日预测处理量



\begin{equation}
B_{k,t}=\sum_i \hat y_{i,k,t}+\sum_j \hat y_{k,j,t}.
\end{equation}



由此得到预测总处理量 $T_k^{f}=\sum_t B_{k,t}$、预测日均值 $\bar B_k$、预测标准差 $\sigma_k$。同理，用历史数据计算历史总处理量 $T_k^{h}$。网络结构中，场地入度和出度分别为



\begin{equation}
d_k^{in}=|\{i:y_{i,k}^{his}>0\}|,\quad d_k^{out}=|\{j:y_{k,j}^{his}>0\}|.
\end{equation}



为消除量纲差异，本文对每个正向指标使用百分位秩归一化。设全体场地数为 $N_s$，指标 $a_k$ 越大越重要，则



\begin{equation}
\rho_s(a_k)=\frac{\operatorname{rank}(a_k)}{N_s},
\end{equation}



其中 $\operatorname{rank}(a_k)$ 按从小到大排序，最大值接近 1。场地重要性从四个分项计算：



\begin{enumerate}

\item 货量规模得分 $V_k$：同时考虑预测期业务量和历史业务量。由于历史数据存在空白、临时异常和长期停运式异常，预测结果更能代表 2023 年 1 月正常基准，因此取

\begin{equation}
V_k=0.75\rho_s(T_k^f)+0.25\rho_s(T_k^h).
\end{equation}

\item 网络结构得分 $G_k$：用场地连接广度衡量其在网络中的枢纽性，定义为

\begin{equation}
G_k=\rho_s(d_k^{in}+d_k^{out}).
\end{equation}

\item 波动性得分 $W_k$：用预测日处理量变异系数衡量调度不稳定性，

\begin{equation}
W_k=\rho_s\left(\frac{\sigma_k}{\bar B_k+\varepsilon}\right),
\end{equation}

其中 $\varepsilon$ 为极小正数，用于避免分母为 0。

\item 关停直接影响得分 $F_k$：若场地 $k$ 关停，直接需要重新分配的货量近似为所有与该场地相关的预测处理量，因此

\begin{equation}
F_k=\rho_s(T_k^f).
\end{equation}

\end{enumerate}



综合得分权重按“业务规模优先、结构与失效影响并重、波动作为辅助”的原则设定。货量规模直接反映场地承担的常态业务，因此给最高权重 0.35；网络结构决定其连接上下游的范围，关停直接影响决定失效后需要重分配的货量，二者同为关键影响项，各取 0.25；波动性主要反映调度难度，重要但不应压过规模和结构，因此取 0.15。四项权重和为 1，且不引入负荷率，避免把短期超载风险误当作结构重要性。于是场地重要性得分为



\begin{equation}
I_k^{site}=0.35V_k+0.25G_k+0.15W_k+0.25F_k,
\end{equation}



其中 $V_k$、$G_k$、$W_k$、$F_k$ 分别为货量规模、网络结构、波动性和关停直接影响得分。



场地重要性前 15 名如下：




\begin{figure}[H]
  \centering
  \includegraphics[width=0.95\textwidth]{\detokenize{figures/table13_site_importance_bar.png}}
  \caption{问题四场地重要性前十五名}
  \label{fig:site-importance-bar}
\end{figure}




可以看到 DC14、DC9、DC10、DC3、DC4、DC8 等场地同时具有较大预测货量和较强网络连接能力，是网络中的核心节点。



\subsection{线路重要性评价模型}



对线路 $e=(i,j)$，记预测总货量 $Q_e^f=\sum_t \hat y_{i,j,t}$，历史总货量 $Q_e^h=\sum_t y_{i,j,t}^{his}$，预测日均值 $\bar q_e$，预测标准差 $\sigma_e$，线路运输能力为 $C_e^{cap}$。同样用百分位秩归一化。若线路总数为 $N_e$，正向指标 $a_e$ 的归一化为



\begin{equation}
\rho_e(a_e)=\frac{\operatorname{rank}(a_e)}{N_e}.
\end{equation}



线路重要性从四个分项计算：



\begin{enumerate}

\item 货量规模得分 $V_e$：预测货量更能代表修正后的正常需求，历史货量用于保留长期业务惯性，因此

\begin{equation}
V_e=0.75\rho_e(Q_e^f)+0.25\rho_e(Q_e^h).
\end{equation}

\item 端点重要性得分 $H_e$：若线路连接两个重要场地，其结构价值更高。先取

\begin{equation}
E_{i,j}=0.5I_i^{site}+0.5I_j^{site},
\end{equation}

再令 $H_e=\rho_e(E_{i,j})$。

\item 波动性得分 $W_e$：用预测日货量变异系数表示线路流量不稳定程度，

\begin{equation}
W_e=\rho_e\left(\frac{\sigma_e}{\bar q_e+\varepsilon}\right).
\end{equation}

\item 替代风险得分 $R_e$：线路 $i\to j$ 的替代能力由同起点其他线路和同终点其他线路的能力共同决定。定义

\begin{equation}
A_e=\sum_{j'\ne j}C_{i,j'}^{cap}+\sum_{i'\ne i}C_{i',j}^{cap},
\end{equation}

则替代能力相对越弱，线路越不可替代，故取

\begin{equation}
R_e=\rho_e\left(-\frac{A_e}{C_e^{cap}+\varepsilon}\right).
\end{equation}

\end{enumerate}



线路综合权重与场地保持一致，以保证两类排名解释口径统一：货量规模 $V_e$ 取 0.35，端点重要性 $H_e$ 和替代风险 $R_e$ 各取 0.25，波动性 $W_e$ 取 0.15。这样既突出线路自身承载的业务量，也保留其连接核心场地和不可替代性的影响。线路重要性得分为



\begin{equation}
I_e^{edge}=0.35V_e+0.25H_e+0.15W_e+0.25R_e,
\end{equation}



其中 $V_e$、$H_e$、$W_e$、$R_e$ 分别为货量规模、端点重要性、波动性和替代风险得分。



线路重要性前 20 名如下：




\begin{figure}[H]
  \centering
  \includegraphics[width=0.95\textwidth]{\detokenize{figures/table14_edge_importance_bar.png}}
  \caption{问题四线路重要性前二十名}
  \label{fig:edge-importance-bar}
\end{figure}




该结果说明，DC9 并不是唯一需要关注的特殊场地，而是高关停影响场地中的代表。与 DC9 类似，DC14、DC10、DC3、DC4、DC8、DC62 等核心场地及其关键上下游，也需要在新增场地设计中一并考虑。



\subsection{新增场地和线路设计}



新增场地记为 DC\_NEW。线路设计遵循以下原则：



\begin{enumerate}

\item 覆盖核心重要场地：优先连接场地重要性排名靠前的节点。

\item 覆盖高关停影响场地的关键上下游：不仅考虑 DC9，也考虑与其他核心节点高度相关的上游和下游。

\item 缓解现有超载路径：对基准预测中已经出现超载的线路，补充必要方向，使货量可经 DC\_NEW 绕行。

\item 不默认全双向：根据场地角色确定方向。若某场地出货占比不低于 0.6，视为 source，设置 $k\to \mathrm{DC\_NEW}$；若入货占比不低于 0.6，视为 sink，设置 $\mathrm{DC\_NEW}\to k$；若兼具中转属性，则设置双向。

\end{enumerate}



最终连接 14 个既有场地：



DC14、DC9、DC10、DC3、DC4、DC8、DC5、DC62、DC25、DC51、DC15、DC36、DC44、DC23。



新开 16 条有向线路：




\begin{table}[H]
  \centering
  \caption{问题四新增 DC\_NEW 有向线路及能力}
  \label{tab:auto-15}
  \begin{tabular*}{\textwidth}{@{}L{0.50\textwidth}@{\extracolsep{\fill}}L{0.42\textwidth}@{}}
    \toprule
新线路 & 能力 \\
\midrule
DC14->DC\_NEW & 380,000 \\
DC\_NEW->DC9 & 355,000 \\
DC\_NEW->DC10 & 166,000 \\
DC\_NEW->DC3 & 151,000 \\
DC\_NEW->DC4 & 178,000 \\
DC\_NEW->DC8 & 186,000 \\
DC\_NEW->DC5 & 100,000 \\
DC\_NEW->DC62 & 138,000 \\
DC25->DC\_NEW & 100,000 \\
DC51->DC\_NEW & 100,000 \\
DC15->DC\_NEW & 100,000 \\
DC36->DC\_NEW & 100,000 \\
DC44->DC\_NEW & 100,000 \\
DC23->DC\_NEW & 100,000 \\
DC9->DC\_NEW & 355,000 \\
DC\_NEW->DC14 & 380,000 \\
    \bottomrule
  \end{tabular*}
\end{table}




新线路能力采用“随连接场地规模调整、但不超过历史可行能力”的规则。对与既有场地 $k$ 相连的新增线路，记 $P_k$ 为场地 $k$ 在 2023 年 1 月预测期内的最大日处理量，即 $P_k=\max_t B_{k,t}$。新增线路能力取



\begin{equation}
C_{new}=\min(505833,\max(100000,0.55P_k)),
\end{equation}



并向上取整到千位。该式包含三层含义：第一，$0.55P_k$ 使连接 DC14、DC9 等核心场地的线路获得更高能力，能够承接其峰值处理量中约一半的缓冲分流；第二，$100000$ 是最低基础能力，保证 DC5、DC25、DC51 等中小规模或风险场地即使 $0.55P_k$ 较低，也有足够的应急分流通道；第三，$505833$ 是历史已有线路日货量能力的最大值，新线路能力不超过该上界，避免设计出超过现有网络能力口径的过大线路。因此该规则既体现“核心场地多配能力”，又控制新线路建设强度。



新增场地处理能力设为



\begin{equation}
S_{new}=800000.
\end{equation}



该值用于约束 DC\_NEW 每日可处理的入货量与出货量之和。设置为 800000 的依据是：按上述 16 条有向线路进行基准分流后，DC\_NEW 的最大日处理量约为 427962 件，对应最大负荷率约为 0.535；在 10\% 和 15\% 随机扰动检验下，平均最大日处理量约为 429878 件和 437095 件，对应负荷率约为 0.537 和 0.546。也就是说，即使考虑预测随机性，DC\_NEW 仍保留约 45\% 的处理能力余量。因此 $S_{new}=800000$ 既能承担主干缓冲场地的分流任务，又不会因为能力过低而在新增结构中形成新的瓶颈。



\subsection{方案合理性说明}



首先，该方案不是针对 DC9 单独补丁式设计，而是从“高关停影响场地及其关键上下游”推广得到。DC9 之所以在结果中突出，是因为其在问题三中被指定关停，且预测恢复后受影响货量达到 13,700,966 件，位列重要场地第 2，相关关键线路在前 20 名中占比较高。但类似 DC9 的高影响节点也存在，例如 DC14、DC10、DC3、DC4、DC8，因此新增线路同时覆盖这些节点。



其次，方案没有采用全双向连接。若 14 个既有场地全部与 DC\_NEW 双向连接，需要 28 条新线路；本文方案只开 16 条，减少 12 条线路，既降低建设和运营成本，也避免大量与实际流向不一致的冗余线路。对于明显 source 型场地，主要设置指向 DC\_NEW 的线路；对于 sink 型场地，主要设置由 DC\_NEW 指向该场地的线路；只有在超载绕行或关键路径需要时才补充反向线路。例如 DC14 和 DC9 虽有明显角色特征，但因关键路径和超载缓解需要，设置了额外方向。



最后，从鲁棒性改善结果看，新增场地能显著降低超载风险：




\begin{table}[H]
  \centering
  \caption{问题四新增场地鲁棒性改善结果}
  \label{tab:auto-16}
  \resizebox{\textwidth}{!}{%
  \begin{tabular}{lrrr}
    \toprule
情形 & 超载记录改善率 & 超载总量改善率 & 负荷率标准差改善率 \\
\midrule
基础情形 & 46.67\% & 58.81\% & 7.17\% \\
10\% 随机扰动 & 42.47\% & 56.43\% & 7.16\% \\
15\% 随机扰动 & 36.13\% & 54.40\% & 7.14\% \\
    \bottomrule
  \end{tabular}}
\end{table}




在基础情形下，线路超载记录由 15 条下降到 8 条，超载总量由 42,203 件下降到 17,385 件；在 10\% 和 15\% 随机扰动下，超载总量改善率仍分别达到 56.43\% 和 54.40\%。这说明方案不仅能改善基准状态，也能在需求波动下保持一定稳定性。




\begin{figure}[H]
  \centering
  \includegraphics[width=0.95\textwidth]{\detokenize{../../q4_analysis/figures/problem4_new_site_robustness_comparison.png}}
  \caption{问题四新增场地方案鲁棒性改善对比}
\end{figure}




\section{模型评价}



本文模型的优点主要包括：



\begin{enumerate}

\item 预测模型采用“异常修正 + 机器学习预测 + 总量校准”的组合结构，既能利用 LightGBM 表达非线性时间特征，又能通过后处理约束修正长期异常停运造成的系统性偏差，预测结果具有较强的稳定性和业务可解释性。

\item 问题二和问题三的分流模型采用分层优化思想，将“正常流转”置于最高优先级，再依次控制变化线路数和负荷均衡性，目标层次清晰，能够避免不同目标简单加权时因权重选择不当导致主目标被削弱。

\item MILP 模型能够显式刻画线路容量、场地处理能力、线路启用状态、货量守恒和变化线路数等约束，解的含义明确，便于从结果表中追踪每一部分货量的去向。

\item 问题三在最终方案选择中引入统一指标对比，将变化线路数和负荷均衡程度放在同一评价口径下比较，使“少改动”和“均衡负荷”的折中关系更加清楚。

\item 问题四将节点重要性评价与新增场地方案设计分开建模，先识别核心节点和关键线路，再进行方向化建线与鲁棒性检验，模型链条完整，结论不依赖单一指标。

\end{enumerate}



模型不足主要包括：



\begin{enumerate}

\item 预测模型虽然具有较好的非线性拟合能力，但 LightGBM 本身不直接表达物流网络的拓扑结构，线路之间的联动关系主要通过特征和后处理体现，网络传播效应刻画仍不够充分。

\item 分层 MILP 的优先级顺序较强，一旦前一层目标固定，后一层目标只能在有限空间内优化，因此模型更适合题目要求明确优先级的场景；若实际运营更强调成本、均衡和改动幅度的连续权衡，则需要改为多目标 Pareto 或加权目标模型。

\item 关停分流模型按日独立求解，能够降低计算规模并提高可操作性，但没有把相邻日期之间的线路开闭连续性、车辆调度衔接和跨日缓冲能力作为状态变量，时间连续决策表达不足。

\item MILP 模型在候选线路较多时求解速度较慢，尤其是问题三允许动态新开线路后，二进制变量数量明显增加，计算效率对求解器和时间限制较敏感。

\item 新增场地方案采用“重要性筛选 + 规则建线 + 鲁棒性仿真”的方法，计算简洁、解释性强，但不是严格的全局最优选址-建线-容量联合优化模型，因此更适合作为可执行的规划方案，而不是理论最优解。

\end{enumerate}



\section{结论}



本文围绕电商物流网络构建了完整的预测与优化模型。问题一通过异常修正 LightGBM 得到 2023 年 1 月全网正常运行预测基准，总货量为 53,200,743 件。问题二在 DC5 关停且不新开线路条件下，通过三阶段 MILP 比较 $\delta=0,2,5$ 三种变化线路放宽方案，最终采用 $\delta=2$ 折中方案，实现 3,190,249 件货量正常分流，未正常流转货量为 23,779 件。问题三在 DC9 关停且允许动态新开线路条件下，比较三类方案后选择最少变化 MILP，正常流转率达到 0.999090，平均每日总变化线路数为 37.161290 条。问题四进一步评价场地和线路重要性，设计新增场地 DC\_NEW 及 16 条有向线路，在基础情形下使超载记录减少 46.67\%，超载总量减少 58.81\%。



综合来看，本文方案能够在保证包裹正常流转的基础上控制网络改动幅度，并通过新增场地和方向化建线提升网络鲁棒性，具有较好的可解释性和可执行性。



\section{参考文献}



[1] Ke G, Meng Q, Finley T, et al. LightGBM: A Highly Efficient Gradient Boosting Decision Tree. Advances in Neural Information Processing Systems, 2017.



[2] Nemhauser G L, Wolsey L A. Integer and Combinatorial Optimization. Wiley, 1988.



[3] Ahuja R K, Magnanti T L, Orlin J B. Network Flows: Theory, Algorithms, and Applications. Prentice Hall, 1993.



\section{附录1}



介绍：支撑材料的文件列表



问题一支撑材料：revised\_forecast\_all\_edges.csv 文件



问题一支撑材料：revised\_forecast\_target\_edges.csv 文件



问题一支撑材料：problem1\_before\_after\_validation\_metrics.csv 文件



问题一支撑材料：daily\_total\_before\_after.png 文件



问题一支撑材料：key\_sites\_before\_after.png 文件



问题一的代码：problem1\_complete.py



问题二支撑材料：problem2\_dc5\_integer\_tradeoff\_summary.csv 文件



问题二支撑材料：problem2\_dc5\_integer\_tradeoff\_daily\_metrics.csv 文件



问题二支撑材料：problem2\_dc5\_integer\_tradeoff\_allocation\_detail.csv 文件



问题二支撑材料：problem2\_scheme\_summary\_compare.png 文件



问题二的代码：problem2\_complete.py



问题三支撑材料：problem3\_balance\_metric\_comparison.csv 文件



问题三支撑材料：problem3\_dc9\_dynamic\_summary.csv 文件



问题三支撑材料：problem3\_dc9\_dynamic\_allocation\_detail.csv 文件



问题三支撑材料：problem3\_balance\_metric\_comparison.png 文件



问题三的代码：problem3\_complete.py



问题四支撑材料：problem4\_site\_core\_importance.csv 文件



问题四支撑材料：problem4\_edge\_core\_importance.csv 文件



问题四支撑材料：problem4\_new\_site\_lines.csv 文件



问题四支撑材料：problem4\_new\_site\_robustness\_summary.csv 文件



问题四支撑材料：problem4\_new\_site\_robustness\_comparison.png 文件



问题四的代码：problem4\_complete.py

\CumcmMainMatterEnd

\end{document}
